\documentclass{beamer}
\usetheme{Rochester}

\usepackage[T1]{fontenc}
\usepackage[utf8]{inputenc}
\usepackage[polish]{babel}

\title{PlayLink}
\author{
	\texorpdfstring{  % bez tego tex sie pulta do \\
		Bartosz Kołaciński 251554 \and
		Mateusz Kosowski 251558 \\
		Nikodem Nowak 251598 \and
		Jakub Rosiak 251620 \\
		Jakub Rusek 247774
	}{}
}
\date{\today}

\begin{document}
	
	\frame{\titlepage}
	
	\begin{frame}
		\frametitle{Cele i zakres projektu}
		
		Celem projektu jest stworzenie aplikacji webowej \textbf{PlayLink}, umożliwiającej użytkownikom szybkie tworzenie i wyszukiwanie sesji do wspólnej gry w niszowe, retro oraz mniej popularne tytuły multiplayer. Zakres projektu obejmuje przygotowanie interfejsu użytkownika, backendu obsługującego sesje i komunikację w czasie rzeczywistym oraz implementację podstawowych funkcji takich jak tworzenie, przeglądanie, dołączanie i opuszczanie sesji, a także czat w obrębie pokoju.
	\end{frame}
	
	\begin{frame}[shrink=10]
		\frametitle{Podział ról i zadań w zespole}
		
		\begin{itemize}
			\item Jakub Rusek – Backend Developer \& DevOps. Odpowiada za architekturę, autoryzację i automatyzację CI/CD w GitHub Actions. Wdraża strategię rolling deployment, dbając o stabilność i dostępność aplikacji.
			\item Nikodem Nowak – Backend Developer. Zajmuje się logiką serwerową, implementacją API oraz komunikacją realtime przez WebSockety. Odpowiada za zarządzanie sesjami i bazą danych.
			\item Bartosz Kołaciński – Frontend Developer. Odpowiada za budowę interfejsu użytkownika (UI), zapewnienie dobrego UX oraz integrację frontendu z backendem.
			\item Mateusz Kosowski – Solution Architect and Backend Developer. Odpowiada za projektowanie architektury rozwiązania, spójność decyzji technicznych oraz koordynację integracji kluczowych komponentów systemu.
			\item Jakub Rosiak – QA \& Documentation and Frontend Developer. Odpowiada za przygotowanie scenariuszy testowych, walidację kluczowych funkcji MVP oraz utrzymanie spójnej dokumentacji projektowej i postępu prac.
		\end{itemize}
	\end{frame}
	
	\frame{\huge{Analiza istniejących rozwiązań}}
	
	\begin{frame}
		\frametitle{Analiza istniejących rozwiązań -- Discord (serwery LFG)}
		
		Zalety:
		\begin{itemize}
			\item Bardzo szybka i bezpośrednia komunikacja
			\item Ogromna, aktywna baza użytkowników
		\end{itemize}
		Wady:
		\begin{itemize}
			\item Brak dedykowanego systemu sesji (pokoje znikają w gąszczu wiadomości)
			\item Wymaga założenia i konfiguracji konta przed użyciem
			\item Trudności z filtrowaniem konkretnych parametrów rozgrywki
		\end{itemize}
	\end{frame}
	
	\begin{frame}
		\frametitle{Analiza istniejących rozwiązań -- Reddit (np. r/lfg)}
		
		Zalety:
		\begin{itemize}
			\item Szybka komunikacja w formie asynchronicznej
			\item Bardzo duża baza użytkowników
		\end{itemize}
		Wady:
		\begin{itemize}
			\item Całkowity brak systemu sesji w czasie rzeczywistym i filtrowania
			\item Wymaga założenia i konfiguracji konta przed użyciem
			\item Brak sztywnej struktury; posty szybko zostają przykryte nowymi treściami
		\end{itemize}
	\end{frame}
	
	\begin{frame}
		\frametitle{Analiza istniejących rozwiązań -- Faceit}
		
		Zalety:
		\begin{itemize}
			\item Spora baza użytkowników skupionych na e-sporcie
			\item Zaawansowany system dobierania graczy na podstawie umiejętności (skill-based)
		\end{itemize}
		Wady:
		\begin{itemize}
			\item Ignorowanie gier niszowych i mniej popularnych tytułów
			\item Wymaga założenia i konfiguracji konta przed użyciem
			\item Bardzo silny nacisk wyłącznie na rozgrywkę kompetetywną (competitive)
		\end{itemize}
	\end{frame}
	
\end{document}